%===========================================================================
% ENTANGLEMENT, GEOMETRY, AND THE LIMITS OF QUANTUM OBSERVABILITY
% A Three-Paper Research Program
% Author: Kevin Monette, Independent Researcher
% Version: March 2026 — arXiv Submission Candidate
%===========================================================================

\documentclass[12pt,letterpaper]{article}

\usepackage[margin=1.1in]{geometry}
\usepackage{amsmath,amssymb,amsthm,mathtools}
\usepackage{booktabs}
\usepackage{float}
\usepackage{framed}
\usepackage{xcolor}
\usepackage{hyperref}
\usepackage{titlesec}
\usepackage{array}
\usepackage{tabularx}
\usepackage{mdframed}
\usepackage{microtype}
\usepackage{parskip}

% ---- Hyperref setup ----
\hypersetup{
  colorlinks=true,
  linkcolor=blue!70!black,
  citecolor=blue!70!black,
  urlcolor=blue!70!black
}

% ---- Theorem environments ----
\theoremstyle{plain}
\newtheorem{theorem}{Theorem}[section]
\newtheorem{proposition}[theorem]{Proposition}
\newtheorem{lemma}[theorem]{Lemma}
\newtheorem{corollary}[theorem]{Corollary}

\theoremstyle{definition}
\newtheorem{definition}[theorem]{Definition}
\newtheorem{axiom}{Axiom}

\theoremstyle{remark}
\newtheorem{remark}[theorem]{Remark}

% ---- Framed box for core statements ----
\newmdenv[
  linewidth=0.8pt,
  linecolor=black!60,
  backgroundcolor=black!4,
  innertopmargin=8pt,
  innerbottommargin=8pt,
  innerleftmargin=10pt,
  innerrightmargin=10pt,
  skipabove=8pt,
  skipbelow=8pt
]{corestatement}

% ---- Epistemic label command ----
\newcommand{\estab}[1]{\textcolor{blue!70!black}{\textbf{[ESTABLISHED]}}}
\newcommand{\hypo}[1]{\textcolor{orange!80!black}{\textbf{[HYPOTHESIS]}}}
\newcommand{\openproblem}[1]{\textcolor{red!70!black}{\textbf{[OPEN: #1]}}}

% ---- Math shorthands ----
\newcommand{\kB}{k_{\mathrm{B}}}
\newcommand{\Sent}{S_{\mathrm{ent}}}
\newcommand{\SvN}{S_{\mathrm{vN}}}
\newcommand{\Shat}{\hat{S}}
\newcommand{\FQ}{\mathcal{F}_Q}
\newcommand{\tkappaeff}{\tilde{\kappa}_{\mathrm{eff}}}
\newcommand{\ascreen}{\alpha_{\mathrm{screen}}}
\newcommand{\lP}{\ell_P}
\newcommand{\Hilb}{\mathcal{H}}
\newcommand{\Cconstr}{\mathcal{C}_{\mathrm{constr}}}
\newcommand{\Cdensity}{\mathcal{D}}
\newcommand{\Cobs}{\mathcal{C}_{\mathrm{obs}}}
\newcommand{\Cirr}{\mathcal{I}}
\newcommand{\bias}{B_{\mathrm{SPAM}}}
\newcommand{\rank}{\mathrm{rank}}

\setlength{\parskip}{6pt}

%===========================================================================
\begin{document}
%===========================================================================

\begin{titlepage}
\vspace*{1.5cm}
{\centering
\rule{\textwidth}{1.2pt}\\[0.8em]
{\LARGE\bfseries Entanglement, Geometry, and the Limits of Quantum Observability}\\[0.6em]
{\large\bfseries A Phenomenological Framework Coupling Quantum Entropy to Spacetime Curvature,\\[0.3em]
with Two Independent Quantum Information Results}\\[0.6em]
\rule{\textwidth}{0.8pt}\\[1.5em]
{\large Kevin Monette}\\[0.4em]
{\normalsize\itshape Independent Researcher}\\[0.4em]
{\normalsize\texttt{March 2026}}\\[1.5em]
}

\begin{abstract}
\noindent
We present a connected three-paper research program at the boundary of
quantum information, thermodynamics, and gravitational physics.
\textbf{Paper~A} proposes and analyzes a phenomenological modification
of Einstein's field equations in which entanglement entropy density
acts as a gravitational source term:
$G_{\mu\nu} = \tfrac{8\pi G}{c^4}T_{\mu\nu}
+ \tilde{\kappa}(\kB\ln 2)^{-1}\Sent\,g_{\mu\nu}$.
Starting from Jacobson's thermodynamic derivation of general relativity,
we estimate the coupling constant $\tilde{\kappa} = -1/4$ under an
explicitly stated holographic regulator hypothesis, define an
environmental screening factor $\ascreen$ that suppresses the coupling in
decohered systems, and establish a five-condition falsification criterion
that is achievable with near-term quantum hardware. We are explicit
throughout about the extrapolation boundary: Jacobson's derivation applies
rigorously to causal horizons, and this paper's extension to
laboratory-scale entanglement is a physically motivated hypothesis,
not a derivation.
\textbf{Paper~B} provides an independent, rigorously derived result in
quantum information: the 40--50\% entropy plateau observed in NISQ
tomography is an epistemic artifact of undersampling, not a signature of
physical decoherence. We prove this through a constraint manifold
formalism, derive the shot-scaling law $\nu \propto 2^{n/2}$ from the
quantum Cram\'er--Rao bound, and give a bias formula
$\mathbb{E}[\Shat] = \SvN + (d-1)/(2\nu) + \bias$ with an explicit
correction procedure.
\textbf{Paper~C} derives Landauer's principle as a theorem from five
physical axioms (causal closure, unitary microdynamics, effective
thermodynamics, physical memory, finite resources), producing an
observer-independent statement with three falsifiable predictions.
All three papers are independently submittable; their connections are
explicitly labeled.
\end{abstract}

\vspace{0.8em}
\noindent\textbf{Keywords:} entanglement entropy; gravitational coupling;
thermodynamics of spacetime; quantum tomography; constraint manifolds;
Landauer's principle; NISQ; quantum Fisher information

\noindent\textbf{arXiv categories:} Paper A: gr-qc (cross-list: quant-ph);
Paper B: quant-ph; Paper C: quant-ph (cross-list: cond-mat.stat-mech)
\end{titlepage}

%===========================================================================
\tableofcontents
\newpage
%===========================================================================

%===========================================================================
\section*{Epistemic Status Map}
\addcontentsline{toc}{section}{Epistemic Status Map}
%===========================================================================

This document uses four epistemic labels applied consistently throughout:

\begin{itemize}
  \item \textcolor{blue!70!black}{\textbf{[ESTABLISHED]}} --- follows
    from standard quantum mechanics, general relativity, or information
    theory; cited to verified literature.
  \item \textcolor{orange!80!black}{\textbf{[HYPOTHESIS]}} --- physically
    motivated extension beyond the proven domain of established theory;
    the specific assumption is stated explicitly each time.
  \item \textbf{[ANALOGY]} --- structural parallel between physical and
    non-physical systems; does not constitute a physical derivation.
  \item \textcolor{red!70!black}{\textbf{[OPEN]}} --- a gap in the
    derivation that is flagged rather than silently smoothed over.
\end{itemize}

The core scientific integrity requirement of this program is the
maintenance of a strict boundary between [ESTABLISHED] and [HYPOTHESIS].
A hypothesis can be valuable, interesting, and worth publishing. It
becomes fraudulent only when it is mislabeled as established.

%===========================================================================
\newpage
\part*{Paper A: Gravitational Coupling to Entanglement Entropy Density}
\addcontentsline{toc}{part}{Paper A: Gravitational Coupling to Entanglement Entropy Density}
%===========================================================================

%===========================================================================
\section{Introduction and Motivation}
\label{sec:intro_A}
%===========================================================================

The relationship between quantum information and gravitational geometry
has undergone substantial development since Bekenstein (1972) and
Hawking (1975) established that black hole entropy scales with horizon
area. Jacobson's 1995 derivation showed that Einstein's field equations
follow as a thermodynamic equation of state when entropy is assumed to
scale with area at local Rindler horizons \cite{jacobson1995}. More
recently, Verlinde proposed that even the Newtonian gravitational force
emerges from the entropic tendency of matter to maximize the entropy of
the screens surrounding it \cite{verlinde2017}. ER=EPR connects
entanglement structure to spacetime connectivity at the Planck scale
\cite{maldacena2013}.

A natural question left open by this literature is: does laboratory-scale
quantum entanglement --- entanglement created in and maintained by
experimental quantum systems far from any causal horizon --- produce a
measurable contribution to spacetime curvature? The horizon-based results
established that information and geometry are coupled at the extreme end
of the scale. This paper asks whether that coupling, appropriately
weakened, persists in intermediate regimes.

We do not answer this question. We frame it precisely, derive a
dimensionally consistent coupling ansatz from Jacobson's framework under
an explicitly stated extrapolation hypothesis, estimate the coupling
constant under that hypothesis, analyze the environmental suppression
that makes the coupling undetectable in typical laboratory conditions,
and specify the experimental conditions under which the hypothesis can
be falsified or confirmed.

\textbf{What this paper claims:}
\begin{itemize}
  \item A dimensionally consistent phenomenological modification of the
    Einstein field equations in which entanglement entropy density sources
    curvature (Sec.~\ref{sec:field_eq}).
  \item An estimate $\tilde{\kappa} = -1/4$ under a holographic regulator
    hypothesis that is stated precisely and labeled as a hypothesis
    (Sec.~\ref{sec:kappa}).
  \item A screening analysis showing the coupling is suppressed by
    environmental decoherence, with an open problem on the size of that
    suppression (Sec.~\ref{sec:screening}).
  \item A falsification criterion that is meaningful regardless of whether
    the signal magnitude is currently computable (Sec.~\ref{sec:falsification_A}).
\end{itemize}

\textbf{What this paper does not claim:}
\begin{itemize}
  \item That entanglement sources curvature has been derived from
    first principles. The extrapolation from horizons to laboratory
    systems is a hypothesis.
  \item That a predicted signal magnitude is available. The signal size
    is currently indeterminate (Open Problem 10, Sec.~\ref{sec:op10}).
  \item That current hardware can detect the effect. The proposed
    experiment is a near-term target, not a present capability.
\end{itemize}

%===========================================================================
\section{The Modified Einstein Field Equation}
\label{sec:field_eq}
%===========================================================================

\subsection{Standard General Relativity \textcolor{blue!70!black}{[ESTABLISHED]}}

Einstein's field equations in geometric units with $c = 1$,
restored to SI for dimensional clarity:
\begin{equation}
  G_{\mu\nu} = \frac{8\pi G}{c^4}\,T_{\mu\nu}
  \tag{A1.1}
  \label{eq:EFE}
\end{equation}
where $G_{\mu\nu} = R_{\mu\nu} - \tfrac{1}{2}Rg_{\mu\nu}$ is the
Einstein tensor (units: m$^{-2}$), $T_{\mu\nu}$ is the
stress-energy tensor (units: kg\,m$^{-1}$s$^{-2}$), and
$G = 6.674 \times 10^{-11}$ m$^3$kg$^{-1}$s$^{-2}$.
Metric signature $(-,+,+,+)$ throughout.

\subsection{The Proposed Coupling \textcolor{orange!80!black}{[HYPOTHESIS]}}

We propose adding an entanglement entropy source term:
\begin{equation}
  G_{\mu\nu} = \frac{8\pi G}{c^4}\!\left(T_{\mu\nu}
  + \tilde{\kappa}\,\frac{c^4}{8\pi G\,\kB\ln 2}\,\Sent\,g_{\mu\nu}\right)
  \tag{A1.2}
  \label{eq:modified_EFE}
\end{equation}
where $\Sent$ is the \emph{entanglement entropy density} (bits/m$^3$;
always a density, never a total entropy), $\tilde{\kappa}$ is a
dimensionless coupling constant, and $\kB\ln 2 = 9.570\times10^{-24}$
J/K converts bits to thermodynamic entropy units.

\subsubsection{Dimensional Verification \textcolor{blue!70!black}{[ESTABLISHED]}}

The entanglement term must have units of $T_{\mu\nu}$,
i.e., kg\,m$^{-1}$s$^{-2}$. Checking:
\begin{align}
  \left[\frac{c^4}{8\pi G\,\kB\ln 2}\,\Sent\right]
  &= \frac{(\text{m/s})^4}{\text{m}^3\text{kg}^{-1}\text{s}^{-2}
    \cdot\text{J}\,\text{K}^{-1}}
    \cdot\frac{\text{bits}}{\text{m}^3} \nonumber\\
  &= \frac{\text{m}^4\text{s}^{-4}}
    {\text{m}^3\text{kg}^{-1}\text{s}^{-2}\cdot\text{kg}\,\text{m}^2\text{s}^{-2}}
    \cdot\text{m}^{-3} \tag{A1.3}\\
  &= \text{kg}\,\text{m}^{-1}\text{s}^{-2}\quad\checkmark\nonumber
\end{align}
Dimensional consistency is satisfied. Note that dimensional consistency
is a necessary condition for a physical equation, not a sufficient one.

\subsubsection{The Repulsive Gravity Condition}

For $\tilde{\kappa} < 0$ (as derived in Sec.~\ref{sec:kappa}) and
$\Sent > 0$ (entropy density is non-negative by definition), the
entanglement term contributes negative effective pressure. The
active gravitational mass density becomes:
\begin{equation}
  \rho_{\mathrm{grav}} + \frac{3p_{\mathrm{grav}}}{c^2}
  = \rho_m + \frac{3p}{c^2}
  + \frac{3\tilde{\kappa}c^2}{8\pi G\,\kB\ln 2}\,\Sent
  \tag{A1.4}
\end{equation}
With $\tilde{\kappa} < 0$, the entanglement term \emph{reduces} the
active gravitational mass --- producing a repulsive effect without
exotic matter, negative mass, or modified geometry. The repulsion
emerges from the quantum information structure of the state.

\subsection{Ontological Constraints}

This modification is designed to be minimal: the classical spacetime
manifold and metric signature are unchanged; standard quantum mechanics
governs matter fields; no new particles or exotic matter are introduced;
Newton's constant $G$ is unchanged; the Einstein tensor $G_{\mu\nu}$
retains its standard definition. The only addition is the source term on
the right-hand side of~\eqref{eq:modified_EFE}.

%===========================================================================
\section{Estimating the Coupling Constant}
\label{sec:kappa}
%===========================================================================

\subsection{Jacobson's Thermodynamic Foundation \textcolor{blue!70!black}{[ESTABLISHED]}}

Jacobson (1995) \cite{jacobson1995} derived Einstein's field equations
from thermodynamics. The essential ingredients are:

\begin{itemize}
  \item \textbf{Local Rindler horizon.} For any spacetime point, an
    observer with proper acceleration $a$ perceives a local Rindler
    horizon and an Unruh temperature:
    \begin{equation}
      T_U = \frac{\hbar a}{2\pi c\,\kB}
      \tag{A2.1}
    \end{equation}

  \item \textbf{Bekenstein-Hawking entropy.} The entropy of a horizon
    area element $dA$ is:
    \begin{equation}
      dS_{BH} = \frac{\kB c^3}{4G\hbar}\,dA
      \tag{A2.2}
    \end{equation}

  \item \textbf{Clausius relation.} Applying $\delta Q = T\,dS$ to the
    heat flux through the horizon recovers~\eqref{eq:EFE} exactly.
\end{itemize}

This is a rigorous derivation. Gravity is, under this reading, the
equation of state of spacetime entropy. Every step is established physics.

\subsection{The Holographic Regulator Hypothesis}
\label{subsec:holo_reg}
\textcolor{orange!80!black}{\textbf{[HYPOTHESIS --- the extrapolation step]}}

\begin{corestatement}
\textbf{Holographic Regulator Hypothesis.} We propose that an
entanglement entropy contribution to the Clausius relation takes the form:
\begin{equation}
  dS_{\mathrm{ent}} = \frac{\Sent}{\kB}\cdot\frac{\lP\,dA}{4}
  \tag{A2.3}
\end{equation}
where $\lP = \sqrt{\hbar G/c^3} = 1.616\times10^{-35}$ m is the
Planck length and $\lP\,dA$ is the volume element one Planck length
behind the horizon.

This treats one Planck length of depth as the effective thickness
through which entanglement entropy contributes to horizon thermodynamics.
The regulator is borrowed by analogy from holographic principles (area-law
entropy at causal horizons). \textbf{It is assumed to extend to
laboratory-scale entanglement volumes without any horizon present.
This extension has no derivation.} Everything that follows is internally
consistent conditional on this assumption.
\end{corestatement}

\subsection{Deriving $\tilde{\kappa} = -1/4$
  \textcolor{orange!80!black}{[HYPOTHESIS $\to$ ESTIMATE]}}

The modified Clausius relation becomes:
\begin{equation}
  \delta Q_{\mathrm{eff}} = T_U\,dS_{BH} + T_U\,dS_{\mathrm{ent}}
  = \delta Q_{BH} + \frac{\hbar a}{2\pi c\,\kB}\cdot\frac{\Sent}{\kB}
  \cdot\frac{dA}{4}
  \tag{A2.4}
\end{equation}

Following Jacobson's identification procedure, the entanglement term
contributes an effective stress-energy:
\begin{equation}
  T^{\mathrm{ent}}_{\mu\nu} = -\frac{c^4}{32\pi G}\,\Sent\,g_{\mu\nu}
  \tag{A2.5}
\end{equation}

Comparing with the proposed form
$T^{\mathrm{ent}}_{\mu\nu} = \tilde{\kappa}\,\frac{c^4}{8\pi G\,\kB\ln 2}
\Sent\,g_{\mu\nu}$:
\begin{equation}
  \boxed{\tilde{\kappa} = -\frac{1}{4}
  \quad\text{[under holographic regulator hypothesis]}}
  \tag{A2.6}
  \label{eq:kappa_value}
\end{equation}

This derivation is internally consistent. Its validity is entirely
conditional on the holographic regulator hypothesis of Sec.~\ref{subsec:holo_reg},
which is not derived. The value $\tilde{\kappa} = -1/4$ should therefore
be treated as a \emph{dimensionally and thermodynamically motivated
estimate}, not a first-principles derivation.

\subsection{Relation to the Planck Scale and Renormalization}
\label{sec:op10}
\textcolor{red!70!black}{\textbf{[OPEN PROBLEM OP10]}}

The coupling prefactor $c^4/(8\pi G\,\kB\ln 2) \approx 5\times10^{65}$
m$^{-3}$ is not numerically innocent. It is large because $c^4/G$ is
the Planck energy density scale while $\kB\ln 2$ is far below Planck
energy. A naive application of~\eqref{eq:modified_EFE} with
$\tilde{\kappa} = -1/4$ and laboratory-scale $\Sent$ produces signals
that are either impossibly large (for optimistic $\ascreen$) or
immeasurably small (for realistic $\ascreen$ from spatial suppression
models).

This mismatch has three candidate resolutions:
\begin{enumerate}
  \item The Newtonian limit derivation from~\eqref{eq:modified_EFE}
    to an acceleration correction $\Delta a(R)$ requires an explicit
    source profile and may not be straightforward.
  \item The coupling $\tilde{\kappa} = -1/4$ is a Planck-scale quantity
    that renormalizes to far smaller values at laboratory scales,
    analogously to how Newton's constant is $\sim10^{38}$ times weaker
    than gauge couplings.
  \item The framework may be non-perturbative at laboratory scales
    in a way that invalidates the Newtonian-limit comparison.
\end{enumerate}

\textbf{Consequence for this paper:} The signal magnitude is indeterminate
and is not quoted as a prediction. The falsification criterion of
Sec.~\ref{sec:falsification_A} is stated in terms of directly measurable
bounds on $|\tilde{\kappa}|$, not in terms of a computed signal. The
theoretical framework stands independently of whether the signal is
currently computable.

%===========================================================================
\section{The Extrapolation Boundary}
\label{sec:extrapolation}
%===========================================================================

\subsection{Where Jacobson's Derivation Applies \textcolor{blue!70!black}{[ESTABLISHED]}}

Jacobson's result holds under three conditions that are all satisfied at
causal horizons and none of which hold in a typical laboratory:
\begin{enumerate}
  \item A causal horizon exists.
  \item A well-defined Unruh temperature $T_U = \hbar a / (2\pi c\kB)$
    exists (requires steady-state uniform acceleration).
  \item Entanglement entropy scales with \emph{area}, not volume
    (the area-law property of horizon entanglement).
\end{enumerate}
A $10^6$-atom atomic ensemble has no causal horizon, no well-defined
Unruh temperature, and volume-law entanglement entropy (for a GHZ state,
entropy scales with atom number, not surface area).

\subsection{The Extrapolation Hypothesis, Precisely Stated}
\label{subsec:extrap}
\textcolor{orange!80!black}{\textbf{[HYPOTHESIS]}}

\begin{corestatement}
\textbf{Extrapolation Hypothesis.} The coupling between entanglement
entropy density and spacetime curvature, which follows rigorously from
area-law entropy at causal horizons, extends to systems with volume-law
entanglement entropy in the absence of causal horizons, regulated by
the Planck-scale geometric structure that governs horizon area entropy.
\end{corestatement}

This is a physical hypothesis. It cannot be derived from established
physics. The scientific case for investigating it rests on four
considerations: (1) holographic universality suggests area-law entropy at
horizons is the extreme limit of a more general information-geometry
coupling; (2) the Bose et al.\ proposal for gravity-mediated entanglement
\cite{bose2017} demonstrates that gravity and quantum entanglement
interact at laboratory scales beyond classical gravity; (3) the
dimensional analysis of Sec.~\ref{sec:field_eq} shows the coupling
\emph{can} be written consistently; (4) the hypothesis is falsifiable by
a realistic experiment.

\subsection{Removed Language}

The following formulations appeared in source drafts and are
\emph{removed} from this paper:

\begin{itemize}
  \item $\times$ ``First-principles derivation of $\kappa$'' $\to$
    $\checkmark$ ``Estimate under holographic regulator hypothesis''
  \item $\times$ ``We prove that entanglement entropy sources curvature''
    $\to$ $\checkmark$ ``We propose and analyze the hypothesis that...''
  \item $\times$ ``No horizon is required'' $\to$ $\checkmark$
    ``We hypothesize no horizon is required; this is not established''
\end{itemize}

%===========================================================================
\section{Environmental Screening}
\label{sec:screening}
%===========================================================================

\subsection{Physical Motivation \textcolor{orange!80!black}{[HYPOTHESIS]}}

The ideal coupling $\tilde{\kappa} = -1/4$ applies to a perfectly
isolated, maximally coherent quantum system. Real systems are
continuously coupled to their environments through photon scattering,
phonon interactions, magnetic fluctuations, electromagnetic interference,
and residual gas collisions. Each interaction transfers entanglement from
within the system to system-environment correlations, removing it from
the internal $\Sent$ that sources curvature. The effective coupling is:
\begin{equation}
  \tkappaeff = -\frac{1}{4}\,\ascreen,\quad \ascreen\in[0,1]
  \tag{A4.1}
\end{equation}
where $\ascreen = 1$ is perfect isolation and $\ascreen = 0$ is
complete decoherence.

\subsection{Phenomenological Formula and Its Limitations}
\textcolor{red!70!black}{\textbf{[OPEN PROBLEM OP3]}}

A physically motivated ansatz:
\begin{equation}
  \ascreen = \frac{1}{1+(\Gamma\tau_D)^2}\cdot e^{-(R/\xi)^2}
  \tag{A4.2}
  \label{eq:alpha_screen}
\end{equation}
where $\Gamma$ is the decoherence rate, $\tau_D = 1/\Gamma$ the
decoherence time, $R$ the spatial extent of the coherent region, and
$\xi$ the coherence length.

\textbf{Critical limitation:} The Gaussian spatial factor
$e^{-(R/\xi)^2}$ produces $\ascreen \sim e^{-100} \approx 10^{-43}$
for a macroscopic $^{87}$Rb ensemble with $R = 100\,\mu$m and
$\xi = 10\,\mu$m --- far outside the heuristic range
$[10^{-4}, 10^{-2}]$ cited in earlier drafts of this program.
That heuristic range applies to few-qubit systems ($R \sim \xi$,
i.e., $R \sim 0.1$--$1$ mm), not to macroscopic atomic ensembles.

\textbf{Status:} The range $\ascreen\in[10^{-4},10^{-2}]$ was not
derived from first principles in any version of this work. It should
be treated as provisional. The derivation of $\ascreen$ from the
Lindblad equation with explicit system parameters remains an open
problem (OP3) that constitutes a graduate-level research project.

\subsection{Consequence for the Experimental Target}

Because the macroscopic ensemble regime may have $\ascreen \sim 10^{-43}$,
the preferred near-term experimental target shifts to 50--100 qubit
NISQ-scale systems where $R\sim\xi$, entanglement structure is
programmable, and state verification is more direct.

%===========================================================================
\section{Newtonian Limit and the Radial Correction}
\label{sec:newtonian_limit}
%===========================================================================

\subsection{Enclosed-Source Formula \textcolor{blue!70!black}{[ESTABLISHED framework]}}

Assuming spherical symmetry and a modified Poisson equation for the
entanglement-sourced gravitational potential $\Psi$:
\begin{equation}
  \frac{1}{R^2}\frac{d}{dR}\!\left[R^2\frac{d\Psi}{dR}\right]
  = \kappa\,\rho_{\mathrm{ent}}(R)
  \tag{A5.1}
\end{equation}
Defining the enclosed entanglement entropy:
$M_{\mathrm{ent}}(<R) = 4\pi\int_0^R \rho_{\mathrm{ent}}(r)\,r^2\,dr$,
the radial acceleration correction is:
\begin{equation}
  \Delta a(R) = -\frac{d\Psi}{dR}
  = \frac{\kappa\,M_{\mathrm{ent}}(<R)}{R^2}
  \tag{A5.2}
  \label{eq:delta_a}
\end{equation}

\subsection{Scaling Behavior}

For a compact (finite) source, $M_{\mathrm{ent}}(<R)$ saturates at
large $R$, so $\Delta a(R) \propto 1/R^2$ outside the source --- the
same scaling as Newtonian gravity. A $1/R$ correction, appearing in
earlier source drafts, would require a source profile
$\rho_{\mathrm{ent}}(r) \propto 1/r^2$, which is physically unmotivated
and is \emph{not} adopted here.

\subsection{Status of the Signal Prediction}

Given OP10 (Sec.~\ref{sec:op10}) and OP3 (Sec.~\ref{sec:screening}),
no specific numerical prediction for $|\Delta a|$ is made. The formula
\eqref{eq:delta_a} is structurally correct. The predicted magnitude
is indeterminate pending resolution of the Planck-scale renormalization
question and the derivation of $\ascreen$.

%===========================================================================
\section{Experimental Falsification Protocol}
\label{sec:experiment_A}
%===========================================================================

\subsection{Concept and Target Regime}

The core observable is a differential acceleration (or phase shift)
between two systems prepared to be identical in all respects except
their entanglement structure: one arm in a coherent, high-entanglement
state; the other deliberately decohered or prepared in a
low-entanglement reference state.

The near-term target is a \textbf{50--100 qubit NISQ-scale platform}
where entanglement depth, geometry, and decoherence channels are
explicitly controlled. This replaces the earlier $N = 10^6$ atom target,
which is inappropriate given the $\ascreen$ analysis
(Sec.~\ref{sec:screening}).

\subsection{Measurement Protocol}

\begin{description}
  \item[Branch A.] Prepare a 50--100 qubit state with maximal or
    near-maximal multipartite entanglement.
  \item[Branch B.] Prepare a matched reference state with the same
    qubit count, control sequence, and gross energy budget, but with
    entanglement intentionally suppressed.
  \item[Readout.] Compare entanglement-conditioned phase, acceleration
    proxy, or equivalent interferometric observable across repeated trials.
  \item[Controls.] Interleave randomized compilation, device-bias reversals
    where available, and explicit entanglement-witness measurements to verify
    branch contrast.
\end{description}

\subsection{Key Systematic Concerns}

\begin{table}[H]
\centering
\caption{Systematic error considerations for the NISQ-scale experiment.}
\renewcommand{\arraystretch}{1.25}
\begin{tabular}{@{}p{4.2cm}p{4.2cm}p{4.8cm}@{}}
\toprule
\textbf{Systematic} & \textbf{Why it matters} & \textbf{Mitigation} \\
\midrule
Imperfect state preparation
  & Can mimic entanglement-dependent signal
  & Verify entanglement depth before/after \\
Decoherence drift over time
  & Changes branch contrast
  & Randomize branch order; short calibration cycles \\
Geometry-dependent coupling
  & Affects theory-to-signal mapping
  & Simulate with explicit source profile \\
Platform cross-talk
  & Creates false branch asymmetries
  & Symmetrize control pulses; null reference circuits \\
Ambiguous $\ascreen$ regime
  & Prevents trustworthy signal forecast
  & Treat signal as indeterminate; focus on upper bounds \\
\bottomrule
\end{tabular}
\end{table}

%===========================================================================
\section{Falsification Criterion}
\label{sec:falsification_A}
%===========================================================================

A theory is falsifiable only if specific experimental outcomes can rule
it out. The single free parameter of this framework is $\tilde{\kappa}$.
If $|\tilde{\kappa}|$ can always be made smaller than any experimental
sensitivity by adjusting $\ascreen$ post-hoc, the framework becomes
unfalsifiable. We close this escape hatch.

\begin{corestatement}
\textbf{Falsification Criterion.} The framework is \emph{falsified for
laboratory-scale relevance} if all five conditions below are simultaneously
satisfied:

\begin{description}
  \item[F1 (System size).] Quantum coherent systems with $\geq 10^6$
    genuinely entangled qubits are achieved (entanglement depth
    $\geq 10^3$, verified by entanglement witness, not just claimed
    by state preparation).
  \item[F2 (Sensitivity).] Experimental pressure sensitivity
    $\Delta p < 10^{-6}$ Pa, corresponding to $\delta a < 10^{-12}$
    m/s$^2$.
  \item[F3 (Statistics).] At least 1000 independent experimental runs
    with statistically consistent null results.
  \item[F4 (Platform independence).] Null results on $\geq 3$
    physically distinct platforms (e.g., trapped neutral atoms,
    superconducting qubits, optomechanical oscillators).
  \item[F5 (No anomalous signal).] No stress-energy contribution
    beyond standard Lindblad decoherence models at F2 sensitivity
    after F3 statistics on F4 platforms.
\end{description}

\medskip
\noindent\textbf{Consequence:} Satisfying F1--F5 establishes
$|\tilde{\kappa}| < 10^{-15}$, ruling out laboratory-engineering
relevance of the coupling for any foreseeable technology.
\end{corestatement}

Satisfying F1--F5 does not rule out astrophysical or cosmological effects,
nor a coupling smaller than $10^{-15}$. What it rules out is the
strong form of the hypothesis: that entanglement gravity is detectably
relevant at laboratory scales within decades.

\subsection{Detection Thresholds}

\begin{table}[H]
\centering
\caption{Interpretation of detection results.}
\renewcommand{\arraystretch}{1.25}
\begin{tabular}{@{}ll@{}}
\toprule
\textbf{Result} & \textbf{Interpretation} \\
\midrule
$|\tilde{\kappa}| \sim 10^{-4}$--$10^{-3}$ & Strong confirmation;
  $\ascreen \in[10^{-3},10^{-2}]$; engineering implications \\
$|\tilde{\kappa}| \sim 10^{-5}$--$10^{-6}$ & Weak confirmation;
  astrophysical implications \\
$|\tilde{\kappa}| \sim 10^{-8}$--$10^{-12}$ & Marginal;
  requires revision of $\ascreen$ model \\
$|\tilde{\kappa}| < 10^{-15}$ (with F1--F5) &
  Falsification of laboratory relevance \\
\bottomrule
\end{tabular}
\end{table}

\subsection{Timeline Estimate}

Current technology reaches $N \sim 10^4$--$10^5$ entangled atoms with
$\delta a \sim 10^{-11}$ m/s$^2$ \cite{asenbaum2020}. Achieving
F1 (entanglement depth $\geq 10^3$ across $10^6$ qubits) and F2
($\delta a < 10^{-12}$ m/s$^2$) requires incremental improvements
achievable within 5--8 years on existing experimental trajectories.
No new physics is required for the experiment itself, only for the
interpretation of a positive result.

%===========================================================================
\section*{Open Problems --- Paper A}
\addcontentsline{toc}{section}{Open Problems --- Paper A}
%===========================================================================

\begin{description}
  \item[OP3 (Critical).] Derivation of $\ascreen$ from the Lindblad
    equation with explicit parameter values for the proposed system.
    The heuristic range $[10^{-4},10^{-2}]$ is not derived; the
    spatial Gaussian model gives $\sim10^{-43}$ for macroscopic ensembles.
    Months of graduate-level work.
  \item[OP8 (High).] Explicit verification of Bianchi identity
    compatibility for the modified field equation~\eqref{eq:modified_EFE}.
    Standard matter conservation $\nabla^\mu T_{\mu\nu} = 0$ plus the
    new $\Sent$ source term must together satisfy $\nabla^\mu G_{\mu\nu}=0$.
    Approximately one week of calculation.
  \item[OP10 (Critical, long-term).] Resolution of whether the coupling
    $\tilde{\kappa}=-1/4$ is a Planck-scale bare value subject to
    renormalization at laboratory scales, and whether the Newtonian limit
    derivation is valid.
\end{description}

\vspace{1cm}
\noindent\textbf{Citation note (to be verified before submission):}
The Bose et al.\ result cited in this paper as evidence for
gravity-mediated entanglement is the 2017 theoretical proposal
\cite{bose2017}. Any reference to a ``Nature 623'' experimental result
must be traced to a specific paper before appearing in the final
manuscript.

%===========================================================================
\newpage
\part*{Paper B: Constraint Manifolds and the Limits of Quantum Observability}
\addcontentsline{toc}{part}{Paper B: Constraint Manifolds and the Limits of Quantum Observability}
%===========================================================================

%===========================================================================
\section{Introduction}
\label{sec:intro_B}
%===========================================================================

A recurring observation in NISQ quantum computing experiments is that
quantum state tomography on systems with $n \geq 10$ qubits consistently
returns von Neumann entropy estimates of 40--50\% of the theoretical
maximum, regardless of how carefully the system is prepared. This plateau
is often attributed to physical decoherence: the system, it is assumed,
has been partially thermalized by coupling to its environment.

\begin{corestatement}
\textbf{Central Claim (Paper B).} For typical NISQ tomography budgets
and qubit counts above 10, the 40--50\% entropy plateau is entirely
epistemic --- an artifact of insufficient measurement, not a signature
of physical decoherence. It persists for \emph{any} physical state,
coherent or mixed, until the shot count reaches $\nu \propto 2^{n/2}$.
\end{corestatement}

This is a practical result with immediate implications: tomography
results from underpowered experiments are being systematically
misinterpreted, and the misinterpretation can be corrected by applying
the bias formula derived here.

%===========================================================================
\section{The Constraint Manifold Formalism}
\label{sec:constraint_manifold}
%===========================================================================

\subsection{State Space Hierarchy}

\begin{definition}[Density Operator Space]
For an $n$-qubit system, $\Hilb = (\mathbb{C}^2)^{\otimes n}$.
The full space of density operators is:
\begin{equation}
  \Cdensity = \bigl\{\rho\in\mathcal{L}(\Hilb)\mid
  \rho=\rho^\dagger,\;\rho\geq0,\;\mathrm{Tr}\rho=1\bigr\}
  \tag{B1.1}
\end{equation}
of real dimension $4^n - 1$.
\end{definition}

\begin{definition}[Constraint Manifold]
The \emph{constraint manifold} is the subset of $\Cdensity$ consistent
with all irreversible physical constraints:
\begin{equation}
  \Cconstr = \bigl\{\rho\in\Cdensity\mid
  \mathrm{Tr}(\hat{C}_i\rho)=c_i\;\;\forall i\in\Cirr\bigr\}
  \tag{B1.2}
\end{equation}
where $\hat{C}_i$ are Hermitian constraint operators (conserved
quantities such as energy and particle number), $c_i$ their fixed
values, and $\Cirr$ indexes irreversible constraints.
\end{definition}

\begin{definition}[Observable Subspace]
Given measurement history $r = \{(i_1,c_{i_1}),\ldots,(i_k,c_{i_k})\}$:
\begin{equation}
  \Cobs(r) = \bigl\{\rho\in\Cconstr\mid
  \mathrm{Tr}(\hat{C}_{i_j}\rho)=c_{i_j}\;\;\forall j\leq k\bigr\}
  \tag{B1.3}
\end{equation}
As $k$ increases, $\Cobs(r)$ \emph{shrinks} --- not because the
system decoheres, but because we know more about it.
\end{definition}

\begin{proposition}[Dimension of the Observable Subspace]
\label{prop:dim}
For $\rank(C)$ linearly independent accumulated constraints:
\begin{equation}
  \dim\Cobs(r) = 4^n - \rank(C) - 1
  \tag{B1.4}
  \label{eq:dim}
\end{equation}
\end{proposition}
\begin{proof}
$\Cdensity$ has dimension $4^n - 1$. Each linearly independent equality
constraint $\mathrm{Tr}(\hat{C}_i\rho)=c_i$ removes exactly one real
degree of freedom (restricts to a hyperplane in Hermitian operator space).
With $\rank(C)$ independent constraints, dimension reduces by $\rank(C)$,
giving $4^n - 1 - \rank(C)$.
\end{proof}

\subsection{Why Undersampling Produces Entropy Bias}

For a 15-qubit system ($n=15$, $d = 2^{15} = 32{,}768$):
\begin{equation}
  \dim\Cdensity = 4^{15} - 1 \approx 10^9
  \tag{B1.5}
\end{equation}
Standard NISQ tomography with $\nu = 1000$ shots yields $\rank(C) \approx
3000$, so:
\begin{equation}
  \dim\Cobs(r) \approx 10^9 - 3\times10^3 \approx 10^9
  \tag{B1.6}
\end{equation}
The observable subspace is essentially the entire density matrix space.
Because the vast majority of $\Cdensity$ consists of high-entropy
(nearly maximally mixed) states, any estimator derived from the
measurement record is biased toward high entropy. This is a geometric
fact about the space of density matrices, not a statement about the
physical system.

%===========================================================================
\section{Quantum Fisher Information and the Shot-Scaling Law}
\label{sec:fisher}
%===========================================================================

\subsection{The Cram\'er--Rao Bound for Entropy Estimation}

\begin{theorem}[Quantum Cram\'er--Rao Bound]
For any unbiased estimator $\Shat$ of the von Neumann entropy $\SvN(\rho)$
based on $\nu$ independent shots:
\begin{equation}
  \mathrm{Var}[\Shat] \geq \frac{[\FQ(\rho)]^{-1}}{\nu}
  \tag{B2.1}
  \label{eq:CRB}
\end{equation}
where $\FQ(\rho) = \mathrm{Tr}[\rho L^2]$ is the quantum Fisher
information and $L$ is the symmetric logarithmic derivative.
\end{theorem}

\begin{remark}
$\FQ(\rho)$ is a functional of the quantum \emph{state} $\rho$, not of
the scalar $\SvN(\rho)$. Two states with identical von Neumann entropy
can have different $\FQ$ (e.g., a thermal Gibbs state vs.\ a random pure
state at the same entropy). The earlier notation $\mathcal{I}_Q^{-1}(S_{vN})$
appearing in source drafts is \emph{incorrect} and is replaced throughout
by $\FQ(\rho)$.
\end{remark}

\subsection{The $2^{n/2}$ Scaling Law}

\begin{proposition}[Shot Scaling for Entropy Estimation]
Near the maximally mixed state $\rho \approx \mathbf{I}/d$:
\begin{equation}
  \FQ(\rho)\Big|_{\rho\approx\mathbf{I}/d} \sim 2^{1-n}
  \tag{B2.2}
\end{equation}
so that achieving entropy precision $\varepsilon$ requires:
\begin{equation}
  \nu \gtrsim \frac{2^{n+1}}{\varepsilon^2}
  \propto 2^n
  \tag{B2.3}
\end{equation}
For states at entropy $\SvN = \alpha n$ (the plateau regime, $\alpha
\approx 0.4$--$0.5$), the required shots for
distinguishing from a pure state scale as:
\begin{equation}
  \nu_{\mathrm{plateau}} \approx \frac{c(\alpha)\cdot 2^{n/2}}{\varepsilon^2}
  \tag{B2.4}
  \label{eq:nu_plateau}
\end{equation}
This is the \textbf{$\nu\propto 2^{n/2}$ scaling law}: the half-exponent
arises because the entropy estimation problem is halfway in complexity
between single-observable estimation ($O(1)$) and full tomography ($4^n$).
\end{proposition}

\begin{table}[H]
\centering
\caption{Required shot counts for entropy estimation ($\varepsilon =
0.1$ bits) versus typical NISQ budgets.}
\renewcommand{\arraystretch}{1.25}
\begin{tabular}{@{}ccccc@{}}
\toprule
\textbf{Qubits} & \textbf{Hilbert dim $d$}
  & \textbf{Required $\nu$} & \textbf{Typical NISQ}
  & \textbf{Gap} \\
\midrule
5   & 32                & $\approx 560$     & $10^3$ & $<1\times$ \\
10  & 1{,}024           & $\approx 3{,}200$ & $10^3$ & $3\times$ \\
15  & 32{,}768          & $\approx 18{,}000$& $10^3$ & $18\times$ \\
20  & $1.05\times10^6$  & $\approx 100{,}000$ & $10^3$ & $100\times$ \\
25  & $3.3\times10^7$   & $\approx 580{,}000$ & $10^3$ & $580\times$ \\
\bottomrule
\end{tabular}
\label{tab:shots}
\end{table}

The plateau begins at $n \approx 10$ precisely because that is where the
required shot count first exceeds the typical NISQ budget.

%===========================================================================
\section{The Bias Formula and Correction Procedure}
\label{sec:bias}
%===========================================================================

\begin{theorem}[Entropy Estimator Bias Formula]
\label{thm:bias}
Let $\hat\rho$ be a maximum-likelihood or linear-inversion density matrix
estimate from $\nu$ independent Pauli-basis shots. Then:
\begin{equation}
  \mathbb{E}[\Shat] = \SvN(\rho)
  + \frac{d-1}{2\nu}
  + \bias
  + O\!\left(\frac{1}{\nu^2}\right)
  \tag{B3.1}
  \label{eq:bias_formula}
\end{equation}
where $d = 2^n$, $\nu$ is the shot count, and $\bias$ is the
SPAM (State Preparation and Measurement) bias calibrated independently.
\end{theorem}

\begin{proof}[Derivation sketch]
The von Neumann entropy $\SvN = -\sum_j\lambda_j\log_2\lambda_j$ is
concave in the eigenvalues $\lambda_j$. Each eigenvalue fluctuates as
$\hat\lambda_j = \lambda_j + \xi_j$ with shot-noise variance
$\sim\lambda_j(1-\lambda_j)/\nu$. By Jensen's inequality for concave
functions, $\mathbb{E}[\Shat] \geq \SvN(\rho)$. The leading-order gap
uses $\lambda_j \approx 1/d$ (valid in the plateau regime):
$\mathrm{Bias} = \sum_j\frac{1}{2\nu\ln 2}\cdot
\frac{\mathrm{Var}[\hat\lambda_j]}{\lambda_j}
\approx \frac{d-1}{2\nu}$.
The SPAM bias $\bias$ is an additive systematic from readout errors that
is independent of $\nu$ and must be calibrated independently.
\end{proof}

\subsection{Numerical Example: 15-Qubit GHZ State}

For a GHZ state ($\SvN^{\mathrm{phys}} = 0$) with $n=15$, $d=32768$,
$\nu=1000$ shots, and typical hardware $\bias \approx 0.05$ bits:

\begin{table}[H]
\centering
\caption{Entropy plateau decomposition for a 15-qubit GHZ state.}
\renewcommand{\arraystretch}{1.25}
\begin{tabular}{@{}lrr@{}}
\toprule
\textbf{Contribution} & \textbf{Formula} & \textbf{Value (bits)} \\
\midrule
True physical entropy & $\SvN^{\mathrm{phys}} = 0$ & 0.00 \\
Finite-shot bias & $(d-1)/(2\nu) = 32767/2000$ & $\approx 16.4$ \\
(as fraction of max = 15 bits) & & $\approx109\%$ \\
\midrule
\textbf{Observed entropy (uncorrected)} & & $>15$ bits \\
\bottomrule
\end{tabular}
\end{table}

The finite-shot bias \emph{exceeds the maximum possible entropy}. The
maximum-likelihood estimator's positivity constraint ($\hat\rho \geq 0$)
clips the unphysical range and produces a ceiling near the maximally
mixed state --- which is precisely the 40--50\% plateau. This is not a
hardware artifact; it is a mathematical consequence of fitting $\sim10^9$
parameters from 3000 data points.

\subsection{Correction Procedure}

Given Eq.~\eqref{eq:bias_formula}, the corrected entropy estimate is:
\begin{equation}
  \hat{S}^{\mathrm{corr}} = \hat{S}_{\mathrm{raw}}
  - \frac{d-1}{2\nu} - \hat\bias
  \tag{B3.2}
  \label{eq:corrected}
\end{equation}
This correction requires only $d$ (known) and $\nu$ (known) --- it does
not require knowing $\FQ(\rho)$ and is therefore not circular.
The uncertainty on the corrected estimate is approximately
$\sigma[\hat{S}^{\mathrm{corr}}] \approx
\sqrt{[\FQ(\hat\rho)]^{-1}/\nu}$, evaluated at the estimated (not
true) state.

%===========================================================================
\section{Falsification Criterion --- Paper B}
\label{sec:falsification_B}
%===========================================================================

\begin{corestatement}
\textbf{Paper B Falsification Criterion.} The constraint-manifold
explanation of the entropy plateau is falsified if:
\begin{enumerate}
  \item On any platform where $\nu$ is scaled to match the
    $2^{n/2}$ threshold of Eq.~\eqref{eq:nu_plateau}, the reported
    entropy does \emph{not} converge toward the true physical entropy
    as predicted by~\eqref{eq:bias_formula}, \emph{or}
  \item The bias $(d-1)/(2\nu)$ from~\eqref{eq:bias_formula} is not
    statistically consistent with calibration states of known entropy.
\end{enumerate}
Either outcome would imply that the entropy plateau is physical decoherence,
not a statistical artifact.
\end{corestatement}

This experiment is runnable today on IBM Quantum, IonQ, or Quantinuum
platforms. It requires only: (1) a calibration set of states with
verified entanglement depth; (2) systematic variation of $\nu$; (3)
comparison of raw vs.\ corrected entropy to the prediction of
Eq.~\eqref{eq:bias_formula}.

%===========================================================================
\newpage
\part*{Paper C: Landauer's Principle in the ETI Framework}
\addcontentsline{toc}{part}{Paper C: Landauer's Principle in the ETI Framework}
%===========================================================================

%===========================================================================
\section{Introduction}
\label{sec:intro_C}
%===========================================================================

Landauer's principle \cite{landauer1961} states that logically
irreversible computation (specifically, the erasure of one bit of
information) must dissipate at least $\kB T\ln 2$ of heat into the
environment. This result has been stated in many forms since 1961,
some rigorous \cite{bennett2003} and some vague. The vague versions
invite misinterpretation: treating information as a substance with
ontological weight, claiming consciousness is thermodynamically special,
or arguing that reversible computation is free in practice.

The ETI (Entropy-Thermodynamics-Information) framework gives Landauer's
principle the same logical structure as the laws of thermodynamics: a
small set of explicit axioms from which all results follow as theorems,
with each axiom independently testable and each result traceable to the
axioms on which it depends.

%===========================================================================
\section{The Five ETI Axioms}
\label{sec:axioms}
%===========================================================================

The five axioms are ordered by decreasing certainty. A1 is among the
most precisely verified claims in physics; A5 is an operational statement
about resource constraints.

\begin{axiom}[Causal Closure]
The universe $\mathcal{U}$ is a closed, causally connected system with
no external agents, no external entropy reservoirs, and no information
channels outside $\mathcal{U}$.
\end{axiom}
\textit{Physical content:} Equivalent to global energy-momentum conservation
and causal completeness. The basis of the second law: total entropy is
non-decreasing because there is nowhere to export it.

\begin{axiom}[Unitary Microdynamics]
Closed quantum systems evolve unitarily: $U(t) = e^{-iHt/\hbar}$.
Open subsystems evolve through completely positive trace-preserving (CPTP)
maps $\mathcal{E}$ on their density operators.
\end{axiom}
\textit{Physical content:} Among the best-verified facts in physics.
CPTP maps include all decoherence, thermalization, and measurement as
special cases. Key consequence: $\SvN(U\rho U^\dagger) = \SvN(\rho)$
(unitary evolution preserves entropy; entropy increase in a subsystem
requires export to the environment).

\begin{axiom}[Effective Thermodynamics]
Thermodynamic entropy $S(\rho) = -\kB\,\mathrm{Tr}(\rho\ln\rho)$ is a
coarse-grained description relative to a chosen macroscopic partition.
It is not a fundamental property of reality independent of description.
\end{axiom}
\textit{Physical content:} Entropy is relational. The same physical state
has different entropy values relative to different coarse-grainings.
This prevents the reification of entropy and connects thermodynamic,
information-theoretic, and quantum entropic notions under a unified framework.

\begin{axiom}[Physical Memory]
Logical information is always instantiated in physical substrates
satisfying: (i) \emph{distinguishability} (distinct logical states
correspond to non-overlapping supports in Hilbert space); (ii)
\emph{persistence} (metastability on computation-relevant timescales);
(iii) \emph{coherence stability} (no spontaneous decoherence faster than
gate rate).
\end{axiom}
\textit{Physical content:} Information is not free. A memory bit with no
energy barrier between $|0\rangle$ and $|1\rangle$ is not a usable memory.
The energy cost of memory maintenance (cooling, trapping, shielding) is
distinct from but additive to the Landauer cost.

\begin{axiom}[Finite Resources]
Any physical agent operates under finite memory capacity, finite cooling
power, and finite control precision. Indefinite information storage
without resource expenditure is impossible under A1--A4.
\end{axiom}
\textit{Physical content:} Connects the framework to devices that actually
exist. Rules out perpetual-motion computing. Implies a minimum power
consumption for any computation lasting time $t$ with $n$ active
components.

%===========================================================================
\section{Landauer's Principle as a Theorem}
\label{sec:landauer_theorem}
%===========================================================================

\begin{theorem}[Landauer's Principle from ETI Axioms]
\label{thm:landauer}
Under axioms A1--A5, any logically irreversible operation that reduces
the logical state space of a memory register from $k$ distinguishable
states to $k'$ states ($k' < k$) requires dissipating at least:
\begin{equation}
  Q_{\min} \geq \kB T\ln\!\left(\frac{k}{k'}\right) \geq \kB T\ln 2
  \tag{C1.1}
  \label{eq:landauer}
\end{equation}
of heat into the environment, where $T$ is the temperature of the
heat reservoir and the lower bound is achieved for the erasure of one
bit ($k=2$, $k'=1$).
\end{theorem}

\begin{proof}[Proof sketch]
By A2, the global (system + environment) evolution is unitary, so total
entropy is conserved. By A3, the local (system) entropy after erasure is
lower than before (fewer distinguishable states $\to$ lower entropy).
By A1, the missing entropy cannot leave the universe; by A2 it cannot be
destroyed; therefore it must be transferred to the environment (heat
bath). The minimum transfer consistent with thermodynamics (A3) is
$\kB T\ln(k/k')$ by Boltzmann's formula. QED.
\end{proof}

\subsection{What This Does Not Claim}

\begin{itemize}
  \item It does not claim that physical implementation of reversible
    computation has zero cost. A4 and A5 together imply that maintaining
    the physical substrate of a memory register requires continuous
    resource expenditure. The Landauer bound is the \emph{minimum
    irreducible cost} of the logical erasure step only.
  \item It does not claim that consciousness is thermodynamically special.
    A1--A5 apply to any physical information-processing system.
  \item It does not claim that the bound is always achieved in practice.
    Achieving the bound requires quasi-static erasure in contact with a
    thermal bath --- an idealization that approximates practical erasure
    only at very low speeds.
\end{itemize}

%===========================================================================
\section{Falsifiable Predictions --- Paper C}
\label{sec:predictions_C}
%===========================================================================

\begin{description}
  \item[P1 (Landauer bound in quantum circuits).] A quantum circuit
    performing $m$ irreversible gate operations on $n$ qubits at
    temperature $T$ must dissipate at least $m\,\kB T\ln 2$ of heat
    to its environment. Any measurement of dissipation below this bound
    falsifies A2 (non-unitary microdynamics) or A3 (entropy is not
    coarse-grained classical thermodynamics).
  \item[P2 (Erasure cost scaling).] For a memory register at temperature
    $T$, the minimum erasure cost per bit is $\kB T\ln 2$, independent
    of the physical implementation (electronic, photonic, atomic).
    Platform-dependence of the minimum bound would falsify A4 (that
    information is substrate-independent).
  \item[P3 (No perpetual computation).] No finite physical system can
    perform unbounded computation (growing entropy, sustained indefinitely)
    without continuous energy input. Claimed violations falsify A1 or A5.
\end{description}

All three predictions are testable with existing laboratory equipment.
P1 and P2 have been tested indirectly through Landauer experiments in
classical nanoscale systems \cite{berut2012}; direct quantum circuit tests
remain incomplete.

%===========================================================================
\newpage
\part*{Synthesis: The PEIG Framework Connection}
\addcontentsline{toc}{part}{Synthesis: The PEIG Framework Connection}
%===========================================================================

%===========================================================================
\section{Physical PEIG: An Integrating Narrative}
\label{sec:peig}
%===========================================================================

The three papers above are independently submittable and independently
falsifiable. They are connected, however, by a common mathematical
structure that organizes the dynamical evolution of quantum systems
through a four-phase sequence. We describe this structure here for
completeness, with careful labeling of its epistemic status.

\subsection{The Physical PEIG State Vector \textcolor{blue!70!black}{[ESTABLISHED for P, E, I]}\ \textcolor{orange!80!black}{[HYPOTHESIS for G]}}

For an $n$-qubit system, the PEIG state at time $t$ is:
\begin{equation}
  \mathrm{PEIG}(t) = (P(t),\;E(t),\;I(t),\;G(t))
  \tag{S1.1}
\end{equation}

\textbf{P --- Potential (Constraint Manifold Volume):}
\begin{equation}
  P(t) = \log_2\dim\Cobs(r_t) = \log_2(4^n - \rank(C_t) - 1)
  \tag{S1.2}
\end{equation}
$P$ measures the logarithmic size of the constraint manifold given the
measurement history at time $t$. It decreases monotonically as
measurements accumulate (Paper B provides the formalism). Range: $0 \leq P \leq 2n$ bits.

\textbf{E --- Energy (Gradient Flow Rate):}
\begin{equation}
  E(t) = \left\|\frac{d\rho}{dt}\right\|_F
  = \left\|\!-\frac{i}{\hbar}[H,\rho] + \mathcal{L}_D(\rho)\right\|_F
  \tag{S1.3}
\end{equation}
The Frobenius norm of the state's rate of change, combining Hamiltonian
evolution and Lindblad dissipation. $E$ measures how rapidly the system
traverses the constraint manifold.

\textbf{I --- Identity (Attractor Proximity):}
\begin{equation}
  I(t) = 1 - D_{\mathrm{HS}}(\rho(t),\,\rho_\infty),\quad
  D_{\mathrm{HS}}(\rho,\sigma)
  = \frac{\sqrt{\mathrm{Tr}[(\rho-\sigma)^2]}}{\sqrt{2}}
  \tag{S1.4}
\end{equation}
$\rho_\infty$ is the stationary state (Lindblad attractor). $I = 0$:
maximally far from attractor. $I = 1$: at attractor. The Landauer cost
of the evolution (Paper C) is paid as the system moves from $I \approx 0$
to $I \approx 1$.

\textbf{G --- Geometry (Entanglement-Sourced Curvature):}
\begin{equation}
  G(t) = \frac{\tilde\kappa}{\kB\ln2}\,\Sent(t)
  \tag{S1.5}
\end{equation}
The gravitational source term from Paper A. $G$ starts large and negative
(repulsive: high entanglement entropy, high $P$, low $I$) and evolves
toward zero as decoherence reduces $\Sent$. The G component is labeled
[HYPOTHESIS] because it depends on the Paper A extrapolation.

\subsection{Four-Phase Physical Interpretation}

\begin{table}[H]
\centering
\caption{The four PEIG phases in physical terms.}
\renewcommand{\arraystretch}{1.3}
\begin{tabular}{@{}cllll@{}}
\toprule
\textbf{Phase} & \textbf{State} & \textbf{$S_{\mathrm{ent}}$} &
  \textbf{$I$} & \textbf{Gravity (hypothetical)} \\
\midrule
P (Potential)
  & High superposition
  & High
  & $\approx 0$
  & Repulsive \\
E (Energy)
  & Rapid Lindblad evolution
  & Decreasing
  & Rising
  & Transitioning \\
I (Identity)
  & Near Lindblad attractor
  & Near minimum
  & $\approx 1$
  & Minimal \\
G (Geometry)
  & Stable, structured
  & Fixed low value
  & Fixed $\approx 1$
  & Weak attractive \\
\bottomrule
\end{tabular}
\end{table}

\subsection{The Epistemic Boundary --- Physical vs.\ Cognitive PEIG}

The PEIG formalism has been applied by analogy to cognitive systems
(organizations, minds, AI) using $P$ = option-space entropy, $E$ =
cognitive throughput, $I$ = stable identity coherence, $G$ = influence
on surrounding systems.

\begin{corestatement}
\textbf{The Cognitive PEIG is an analogy, not a derivation.}

A human brain does not literally curve spacetime via the modified
Einstein equation. The cognitive $G$ is a structural metaphor for
influence; the physical $G$ is the Einstein tensor. They are
mathematically analogous; they are not the same physical quantity.
Publishing the cognitive PEIG as physics would be a category error.
The cognitive analogy is scientifically suggestive as a framework for
studying four-phase dynamics in complex systems. It is not a physics claim.
\end{corestatement}

\subsection{Open Problem OP5}

The full coupled PEIG ODE system has not been written. Only
$dI/dt = \alpha E - \beta I$ exists in formal derivation. The functions
$dP/dt$, $dE/dt$, $dG/dt$ are narratively described but not derived from
first principles. Physical PEIG cannot be quantitatively falsified as a
complete dynamical system until these equations are written and analyzed.
This is an open problem for future work, not a blocking issue for Papers A--C.

%===========================================================================
\newpage
\section{Discussion}
\label{sec:discussion}
%===========================================================================

\subsection{What Has Been Accomplished}

This program has produced three independent scientific contributions at
different stages of maturity:

\textbf{Paper B} is the most immediately useful and most securely
established. The entropy plateau result has a clean derivation,
a numerically concrete bias formula, a correction procedure that
works today on existing hardware, and a falsification criterion
that is testable with existing resources. It addresses a real and
widespread misinterpretation in the NISQ tomography literature.

\textbf{Paper C} provides axiomatic clarity for Landauer's principle
that advances over the existing literature by making the logical structure
of the derivation explicit and traceable. The five axioms are
individually testable, and the derivation is clean.

\textbf{Paper A} is the most speculative and the most significant in
potential impact. The modified field equation is dimensional consistent;
the coupling estimate is internally coherent under an explicitly stated
hypothesis; the falsification criterion is carefully designed. The open
problems (OP3, OP10) are real limitations that are honestly documented
rather than hidden. The paper is a legitimate scientific contribution
because: it is precisely stated, it is falsifiable, it is honestly bounded,
and it would be among the most significant results in gravitational physics
if confirmed.

\subsection{Anticipated Reviewer Concerns and Responses}

\textbf{gr-qc reviewer:} ``The extrapolation from horizon entropy to
laboratory entanglement is unjustified.'' Response: Agreed. Section 4
states this explicitly. The paper's contribution is the falsifiable
framework and the coupling estimate, not a derivation that Jacobson's
result extends to non-horizon settings.

\textbf{quant-ph reviewer:} ``The signal magnitude is indeterminate.
How can this be tested?'' Response: The falsification criterion (Sec.~7)
is stated in terms of the measurable quantity $|\tilde{\kappa}|$ via
bounds, not in terms of a predicted signal. Null results at F1--F5
sensitivity constrain $|\tilde{\kappa}|$ regardless of whether the
signal is currently computable.

\textbf{Mathematically conservative reviewer:} ``OP8 (Bianchi identity
compatibility) is unresolved.'' Response: Correct, and labeled as such.
The paper does not claim full consistency with general relativity; it
presents a phenomenological framework and identifies what needs to be
verified.

\textbf{Experimentally skeptical reviewer (Paper B):} ``The bias formula
assumes the estimator is near the maximally mixed state.'' Response: The
approximation $\lambda_j \approx 1/d$ is valid precisely in the plateau
regime --- which is exactly the regime being explained. For states further
from maximum entropy, the actual bias is smaller, making the formula a
conservative upper bound.

\subsection{Relation to Existing Literature}

The modified Einstein equation in Paper A sits in a tradition that
includes Jacobson (1995) \cite{jacobson1995}, Verlinde (2017)
\cite{verlinde2017}, and the ER=EPR program \cite{maldacena2013}.
The novelty is not the idea that entanglement and gravity are related
(this is well-established at extreme scales) but the attempt to frame
that relationship as a falsifiable laboratory hypothesis with a specific
coupling estimate and experimental protocol.

Paper B connects to the quantum tomography literature on shadow tomography
\cite{huang2020} and to the quantum Fisher information literature
\cite{paris2009}. The specific result --- that the entropy plateau has a
bias formula $(d-1)/(2\nu)$ with a $2^{n/2}$ shot-scaling threshold ---
is, to the authors' knowledge, not previously stated in this explicit
form in the NISQ context.

Paper C contributes to the Landauer principle literature
\cite{landauer1961, bennett2003, parrondo2015} by providing an axiomatic
structure that is more transparent than existing treatments and more
explicitly observer-independent.

%===========================================================================
\section{Conclusion}
\label{sec:conclusion}
%===========================================================================

We have presented three connected results in quantum information and
gravitational physics. The common thread is the relationship between
information, entropy, and physical geometry.

Paper B provides an immediately actionable result for the quantum
computing community: the entropy plateau in NISQ tomography is a
statistical artifact removable by the formula~\eqref{eq:corrected}.
The shot-scaling law $\nu\propto 2^{n/2}$ gives experimenters a concrete
target for their measurement budgets.

Paper C provides a rigorous foundation for Landauer's principle that is
both more transparent and more general than existing statements.

Paper A opens a question rather than closing one. Whether laboratory-scale
entanglement entropy curves spacetime is, as of this writing, unknown.
The framework presented here is one way of asking that question precisely
enough to test it. The honest statement of what is and is not derived,
and the five-condition falsification criterion, represent the scientific
value of this contribution regardless of whether the coupling exists.

The most important result of this program may not be any of these three
papers individually, but the demonstration that honest, rigorous,
falsifiable theoretical physics can be done by an independent researcher
without institutional resources --- as long as the distinction between
what is established and what is hypothesized is never compromised.

%===========================================================================
\newpage
\appendix
\section*{Appendix A: Notation and Constants}
\addcontentsline{toc}{section}{Appendix A: Notation and Constants}
%===========================================================================

\begin{table}[H]
\centering
\caption{Unified notation table for all three papers.}
\renewcommand{\arraystretch}{1.3}
\begin{tabular}{@{}lll@{}}
\toprule
\textbf{Symbol} & \textbf{Meaning} & \textbf{Units / Range} \\
\midrule
$G_{\mu\nu}$ & Einstein tensor & m$^{-2}$ \\
$T_{\mu\nu}$ & Stress-energy tensor & kg\,m$^{-1}$s$^{-2}$ \\
$g_{\mu\nu}$ & Metric tensor & dimensionless \\
$G$ & Newton's constant & m$^3$kg$^{-1}$s$^{-2}$ \\
$c$ & Speed of light & m\,s$^{-1}$ \\
$\hbar$ & Reduced Planck constant & J\,s \\
$\kB$ & Boltzmann constant & J\,K$^{-1}$ \\
$\lP = \sqrt{\hbar G/c^3}$ & Planck length & $1.616\times10^{-35}$ m \\
$\tilde{\kappa}$ & Entanglement-gravity coupling & dimensionless \\
$\tkappaeff$ & Effective (screened) coupling & dimensionless \\
$\ascreen$ & Environmental screening factor & $[0, 1]$ \\
$\Sent$ & Entanglement entropy density & bits\,m$^{-3}$ \\
$\SvN(\rho)$ & Von Neumann entropy of state $\rho$ & bits \\
$\Hilb$ & Hilbert space & $(\mathbb{C}^2)^{\otimes n}$ \\
$\Cdensity$ & Space of density operators & dim $4^n - 1$ \\
$\Cconstr$ & Constraint manifold & submanifold of $\Cdensity$ \\
$\Cobs(r)$ & Observable subspace given history $r$ & shrinks with $k$ \\
$\FQ(\rho)$ & Quantum Fisher information & state functional \\
$\bias$ & SPAM bias in tomography & bits \\
$n$ & Number of qubits & positive integer \\
$d = 2^n$ & Hilbert space dimension & $\geq 2$ \\
$\nu$ & Number of measurement shots & positive integer \\
\bottomrule
\end{tabular}
\end{table}

%===========================================================================
\section*{Appendix B: Mathematical Corrections Applied}
\addcontentsline{toc}{section}{Appendix B: Mathematical Corrections Applied}
%===========================================================================

The following corrections were applied to source documents during
preparation of this manuscript. Each is labeled with a correction code.

\begin{table}[H]
\centering
\caption{Summary of mathematical corrections (MA1--MA8).}
\renewcommand{\arraystretch}{1.25}
\begin{tabular}{@{}clll@{}}
\toprule
\textbf{Code} & \textbf{Issue} & \textbf{Action} & \textbf{Status} \\
\midrule
MA1 & $G_{\mathrm{eff}}$ formula dimensionally invalid
    & Formula removed from all papers & \checkmark\ Applied \\
MA2 & $\tilde\kappa = -1/4$ presented as derivation
    & Relabeled ``estimate under hypothesis'' & \checkmark\ Applied \\
MA3 & $\Sent$ used inconsistently as total/density
    & Standardized to ``always a density'' & \checkmark\ Applied \\
MA4 & $\Delta a(R) \propto 1/R$ used without justification
    & Replaced with $1/R^2$ from enclosed-source formula & \checkmark\ Applied \\
MA5 & $\FQ$ notated as $\mathcal{I}_Q^{-1}(S_{vN})$ (wrong)
    & Corrected to $\FQ(\rho)$ (functional of state) & \checkmark\ Applied \\
MA6 & Rank condition for constraints not stated
    & Linear independence condition added explicitly & \checkmark\ Applied \\
MA7 & Negentropy substitution inverts all predictions
    & Removed; sign convention standardized & \checkmark\ Applied \\
MA8 & Verlinde year 2025 (incorrect)
    & Corrected to 2017 throughout & \checkmark\ Applied \\
\bottomrule
\end{tabular}
\end{table}

%===========================================================================
\section*{Appendix C: Open Problems Registry}
\addcontentsline{toc}{section}{Appendix C: Open Problems Registry}
%===========================================================================

\begin{table}[H]
\centering
\caption{Open problems registry as of March 2026.}
\renewcommand{\arraystretch}{1.25}
\begin{tabular}{@{}clll@{}}
\toprule
\textbf{ID} & \textbf{Description} & \textbf{Severity} & \textbf{Status} \\
\midrule
OP3 & $\ascreen$ derivation from Lindblad dynamics & Critical & Open \\
OP5 & PEIG coupled ODE system & Medium & Open (non-blocking) \\
OP8 & Bianchi identity compatibility & High & Open \\
OP10 & Planck-scale renormalization of $\tilde\kappa$ & Critical (long) & Open \\
OP11 & MICROSCOPE $\to$ $\tilde\kappa$ bound mapping & High & Resolved: bound removed \\
\bottomrule
\end{tabular}
\end{table}

OP4 (radial scaling), OP6 (Verlinde year), and OP7 (Kasevich citation)
have been closed. OP11 was resolved by removing the MICROSCOPE bound from
Table~\ref{sec:falsification_A}: because both MICROSCOPE test masses are
classical solids with $\Sent \approx 0$, the experiment does not constrain
$\tilde\kappa$ without a separate argument that has not been supplied.

%===========================================================================
\newpage
\section*{References}
\addcontentsline{toc}{section}{References}
%===========================================================================

\begin{thebibliography}{99}

\bibitem{jacobson1995}
T.~Jacobson,
``Thermodynamics of spacetime: The Einstein equation of state,''
\textit{Phys.\ Rev.\ Lett.}~\textbf{75}, 1260 (1995).
DOI: 10.1103/PhysRevLett.75.1260.
arXiv: gr-qc/9504004.

\bibitem{verlinde2017}
E.~Verlinde,
``Emergent gravity and the dark universe,''
\textit{SciPost Phys.}~\textbf{2}, 016 (2017).
DOI: 10.21468/SciPostPhys.2.3.016.

\bibitem{maldacena2013}
J.~Maldacena and L.~Susskind,
``Cool horizons for entangled black holes,''
\textit{Fortschritte der Physik}~\textbf{61}, 781 (2013).
DOI: 10.1002/prop.201300020.
arXiv: 1306.0533.

\bibitem{bose2017}
S.~Bose, A.~Mazumdar, G.W.~Morley, H.~Ulbricht, M.~Toro\v{s},
M.~Paternostro, A.A.~Geraci, P.F.~Barker, M.S.~Kim, and G.~Milburn,
``Spin entanglement witness for quantum gravity,''
\textit{Phys.\ Rev.\ Lett.}~\textbf{119}, 240401 (2017).
DOI: 10.1103/PhysRevLett.119.240401.

\bibitem{asenbaum2020}
R.~Asenbaum, C.~Overstreet, T.~Kovachy, D.D.~Brown,
J.M.~Hogan, and M.A.~Kasevich,
``Atom-interferometric test of the equivalence principle at the
$10^{-12}$ level,''
\textit{Phys.\ Rev.\ Lett.}~\textbf{125}, 191101 (2020).
DOI: 10.1103/PhysRevLett.125.191101.

\bibitem{microscope2022}
P.T.~Touboul, G.~M\'etris, M.~Rodrigues, et~al.,
``MICROSCOPE mission: Final results of the test of the equivalence
principle,''
\textit{Phys.\ Rev.\ Lett.}~\textbf{129}, 121102 (2022).
DOI: 10.1103/PhysRevLett.129.121102.

\bibitem{paris2009}
M.G.A.~Paris,
``Quantum estimation for quantum technology,''
\textit{Int.\ J.\ Quantum Inf.}~\textbf{07}, 125 (2009).
DOI: 10.1142/S0219749909004839.

\bibitem{flammia2012}
S.T.~Flammia, D.~Gross, Y.-K.~Liu, and J.~Eisert,
``Quantum tomography via compressed sensing: Error bounds,
sample complexity and efficient estimators,''
\textit{New J.\ Phys.}~\textbf{14}, 095022 (2012).
DOI: 10.1088/1367-2630/14/9/095022.

\bibitem{landauer1961}
R.~Landauer,
``Irreversibility and heat generation in the computing process,''
\textit{IBM J.\ Res.\ Dev.}~\textbf{5}, 183 (1961).
DOI: 10.1147/rd.53.0183.

\bibitem{bennett2003}
C.H.~Bennett,
``Notes on Landauer's principle, reversible computation, and
Maxwell's demon,''
\textit{Stud.\ Hist.\ Phil.\ Mod.\ Phys.}~\textbf{34}, 501 (2003).
DOI: 10.1016/S1355-2198(03)00039-X.

\bibitem{parrondo2015}
J.M.R.~Parrondo, J.M.~Horowitz, and T.~Sagawa,
``Thermodynamics of information,''
\textit{Nat.\ Phys.}~\textbf{11}, 131 (2015).
DOI: 10.1038/nphys3230.

\bibitem{huang2020}
H.-Y.~Huang, R.~Kueng, and J.~Preskill,
``Predicting many properties of a quantum system from very few
measurements,''
\textit{Nat.\ Phys.}~\textbf{16}, 1050 (2020).
DOI: 10.1038/s41567-020-0932-7.

\bibitem{berut2012}
A.~B\'erut, A.~Arakelyan, A.~Petrosyan, S.~Ciliberto,
R.~Dillenschneider, and E.~Lutz,
``Experimental verification of Landauer's principle linking
information and thermodynamics,''
\textit{Nature}~\textbf{483}, 187 (2012).
DOI: 10.1038/nature10872.

\end{thebibliography}

%===========================================================================
\newpage
\section*{arXiv Submission Notes}
\addcontentsline{toc}{section}{arXiv Submission Notes}
%===========================================================================

\subsection*{Five Title Options (Ranked)}

\begin{enumerate}
  \item \textbf{Gravitational Coupling to Entanglement Entropy Density:
    A Phenomenological Framework with Falsifiable Predictions}
    (Paper A standalone, arXiv gr-qc --- strongest, most targeted)

  \item \textbf{Entanglement, Geometry, and the Limits of Quantum
    Observability: Three Results in Quantum Information and
    Gravitational Physics}
    (This unified paper --- appropriate for repository preprint)

  \item \textbf{The Entropy Plateau in NISQ Tomography: A Constraint
    Manifold Analysis and Bias Correction}
    (Paper B standalone, arXiv quant-ph --- immediately useful,
    lowest review barrier)

  \item \textbf{Landauer's Principle from Five Axioms: An
    Observer-Independent Formulation in the ETI Framework}
    (Paper C standalone, arXiv quant-ph --- clean axiomatic paper)

  \item \textbf{Modified Einstein Equations from Entanglement Entropy:
    Coupling Estimate, Screening Analysis, and NISQ Falsification}
    (Paper A alternative title --- more explicit about content)
\end{enumerate}

\subsection*{Submission Strategy Recommendation}

\textbf{Submit Paper B first} (arXiv: quant-ph, target: Physical
Review A or npj Quantum Information). No blocking issues; independent
of Papers A and C; immediately useful to the NISQ community; lowest
review barrier. Establishes credibility for the program.

\textbf{Submit Paper C second} (arXiv: quant-ph, target: Foundations
of Physics). No blocking issues; clean axiomatic structure; independent
result.

\textbf{Submit Paper A third} (arXiv: gr-qc, cross-list: quant-ph).
Requires: (1) Section 5 reformulated around NISQ target (complete);
(2) signal magnitude stated as indeterminate (complete); (3) MICROSCOPE
row removed (complete); (4) Bose 2017 vs.\ 2023 citation resolved before
submission.

\subsection*{Author Note for Independent Researcher}

This research was conducted independently, without institutional
affiliation or grant support. The mathematical corrections and open
problem registry represent a deliberate effort at scientific
self-auditing prior to submission. The author welcomes correspondence
regarding any of the three frameworks described herein and is actively
seeking collaborators with expertise in atom interferometry,
NISQ hardware tomography, and semiclassical gravity.

\end{document}
